\documentclass[a4paper, 14pt]{extarticle}

% Поля
%--------------------------------------
\usepackage{geometry}
\geometry{a4paper,tmargin=2cm,bmargin=2cm,lmargin=3cm,rmargin=1cm}
%--------------------------------------


%Russian-specific packages
%--------------------------------------
\usepackage[T2A]{fontenc}
\usepackage[utf8]{inputenc}
\usepackage[english, main=russian]{babel}

%--------------------------------------

\usepackage{textcomp}

% Красная строка
%--------------------------------------
\usepackage{indentfirst}               
%--------------------------------------             


%Graphics
%--------------------------------------
\usepackage{graphicx}
\graphicspath{ {./images/} }
\usepackage{float}%"Плавающие" картинки
%--------------------------------------

% Полуторный интервал
%--------------------------------------
\linespread{1.3}                    
%--------------------------------------

%Выравнивание и переносы
%--------------------------------------
% Избавляемся от переполнений
\sloppy
% Запрещаем разрыв страницы после первой строки абзаца
\clubpenalty=10000
% Запрещаем разрыв страницы после последней строки абзаца
\widowpenalty=10000
%--------------------------------------

%Списки
\usepackage{enumitem}

\usepackage{minted}
%Подписи
\usepackage{caption} 

%Гиперссылки
% \usepackage{hyperref}

\usepackage[colorlinks, urlcolor = blue, filecolor = blue, citecolor = blue, linkcolor = black]{hyperref}


\hypersetup {
unicode=true
}

%Рисунки
%--------------------------------------
\DeclareCaptionLabelSeparator*{emdash}{~--- }
\captionsetup[figure]{labelsep=emdash,font=onehalfspacing,position=bottom}
%--------------------------------------

\usepackage{tempora}

%Листинги
%--------------------------------------
%Листинги
%--------------------------------------
\usepackage{listings}
\lstset{
  basicstyle=\ttfamily\footnotesize, 
  %basicstyle=\footnotesize\AnkaCoder,        % the size of the fonts that are used for the code
  breakatwhitespace=false,         % sets if automatic breaks shoulbd only happen at whitespace
  breaklines=true,                 % sets automatic line breaking
  captionpos=t,                    % sets the caption-position to bottom
  inputencoding=utf8,
  frame=single,                    % adds a frame around the code
  keepspaces=true,                 % keeps spaces in text, useful for keeping indentation of code (possibly needs columns=flexible)
  keywordstyle=\bf,       % keyword style
  numbers=left,                    % where to put the line-numbers; possible values are (none, left, right)
  numbersep=5pt,                   % how far the line-numbers are from the code
  xleftmargin=25pt,
  xrightmargin=25pt,
  showspaces=false,                % show spaces everywhere adding particular underscores; it overrides 'showstringspaces'
  showstringspaces=false,          % underline spaces within strings only
  showtabs=false,                  % show tabs within strings adding particular underscores
  stepnumber=1,                    % the step between two line-numbers. If it's 1, each line will be numbered
  tabsize=2,                       % sets default tabsize to 8 spaces
  title=\lstname                   % show the filename of files included with \lstinputlisting; also try caption instead of title
}
%--------------------------------------

%%% Математические пакеты %%%
%--------------------------------------
\usepackage{amsthm,amsfonts,amsmath,amssymb,amscd}  % Математические дополнения от AMS
\usepackage{mathtools}                              % Добавляет окружение multlined
\usepackage[perpage]{footmisc}


% графики
\usepackage{booktabs}        % Пакет для улучшенных таблиц
\usepackage{pgfplots}        % Пакет для построения графиков
\pgfplotsset{compat=1.17}    % Установка совместимости

% Литература 

% переименовываем  список литературы в "список используемой литературы"
\addto\captionsrussian{\def\refname{Список литературы}}



\captionsetup[figure]{name=Рисунок, labelsep=space}


%--------------------------------------

%--------------------------------------
%			НАЧАЛО ДОКУМЕНТА
%--------------------------------------

\begin{document}

%--------------------------------------
%			ТИТУЛЬНЫЙ ЛИСТ
%--------------------------------------
\begin{titlepage}
\thispagestyle{empty}
\newpage


%Шапка титульного листа
%--------------------------------------
\vspace*{-60pt}
\hspace{-65pt}
\begin{minipage}{0.3\textwidth}
\hspace*{-20pt}\centering
\includegraphics[width=\textwidth]{emblem}
\end{minipage}
\begin{minipage}{0.67\textwidth}\small \textbf{
\vspace*{-0.7ex}
\hspace*{-6pt}\centerline{Министерство науки и высшего образования Российской Федерации}
\vspace*{-0.7ex}
\centerline{Федеральное государственное автономное образовательное учреждение }
\vspace*{-0.7ex}
\centerline{высшего образования}
\vspace*{-0.7ex}
\centerline{<<Московский государственный технический университет}
\vspace*{-0.7ex}
\centerline{имени Н.Э. Баумана}
\vspace*{-0.7ex}
\centerline{(национальный исследовательский университет)>>}
\vspace*{-0.7ex}
\centerline{(МГТУ им. Н.Э. Баумана)}}
\end{minipage}
%--------------------------------------

%Полосы
%--------------------------------------
\vspace{-25pt}
\hspace{-35pt}\rule{\textwidth}{2.3pt}

\vspace*{-20.3pt}
\hspace{-35pt}\rule{\textwidth}{0.4pt}
%--------------------------------------

\vspace{1.5ex}
\hspace{-35pt} \noindent \small ФАКУЛЬТЕТ\hspace{80pt} <<Информатика и системы управления>>

\vspace*{-16pt}
\hspace{47pt}\rule{0.83\textwidth}{0.4pt}

\vspace{0.5ex}
\hspace{-35pt} \noindent \small КАФЕДРА\hspace{50pt} <<Теоретическая информатика и компьютерные технологии>>

\vspace*{-16pt}
\hspace{30pt}\rule{0.866\textwidth}{0.4pt}

\vspace{11em}

\begin{center}
\Large {\bf Лабораторная работа} \\ 
\large <<Математическая модель лифта>> \\
\large {по курсу <<Моделирование>>} \\

\end{center}\normalsize

\vspace{8em}

\vbox{%
\hfill%
\vbox{%
\hbox{Студент: Аринин М.А.}%
\hbox{Группа: ИУ9-82Б}%
\hbox{Преподаватель: Домрачева А.Б.}%
}%
} 

\bigskip

\vfill


\begin{center}
\textsl{Москва 2026}
\end{center}
\end{titlepage}
%--------------------------------------
%		КОНЕЦ ТИТУЛЬНОГО ЛИСТА
%--------------------------------------


\renewcommand{\ttdefault}{pcr}

\setlength{\tabcolsep}{3pt}

\newpage
\setcounter{page}{2}

\section{Постановка задачи}
Рассматривается задача о движении кабины лифта в вертикальном направлении. Кабина лифта массой $m$ подвешена на тросе и движется под действием силы натяжения троса $T(t)$, силы тяжести $mg$ и силы трения $F_f$, пропорциональной скорости движения.

Лифт работает в различных режимах: разгон при подъеме, равномерное движение, торможение перед остановкой, спуск. Необходимо описать динамику движения кабины при переключении между этими режимами.

Цель работы --- построение математической модели движения лифта, формализация процесса движения через систему дифференциальных уравнений, постановка и решение соответствующих вычислительных задач, а также проверка качества модели.

%--------------------------------------------------------------------
\section{Построение математической модели}
Построение математической модели представляет собой формализацию процесса движения лифта на языке математических формул. Начинаем с формализации системы дифференциальных уравнений, описывающих динамику движения кабины.

Рассматривается движение кабины лифта в вертикальном направлении при следующих допущениях:
\begin{itemize}
  \item кабина лифта рассматривается как материальная точка массой $m$;
  \item трос считается нерастяжимым и невесомым;
  \item сила трения пропорциональна скорости движения: $F_f = -\mu v$;
  \item ускорение свободного падения $g$ постоянно;
  \item движение происходит только в вертикальном направлении.
\end{itemize}

Введем переменные состояния: высота $h(t)$ и скорость $v(t)$ кабины лифта. Эти переменные полностью определяют состояние системы в каждый момент времени.

\subsection{Модель движения лифта}

На кабину лифта действуют следующие силы:
\begin{itemize}
  \item Сила натяжения троса $T(t)$, направленная вверх
  \item Сила тяжести $mg$, направленная вниз
  \item Сила трения $F_f = -\mu v$, направленная против скорости движения
\end{itemize}

Согласно второму закону Ньютона, уравнение движения имеет вид:
\[
m\frac{d^2h}{dt^2} = T(t) - mg - \mu v,
\]
где $v = \frac{dh}{dt}$ --- скорость движения кабины.

Введя переменную состояния $v = \frac{dh}{dt}$, получим систему дифференциальных уравнений первого порядка:
\[
\begin{cases}
  \dfrac{dh}{dt} = v,\\
  \dfrac{dv}{dt} = \dfrac{T(t) - mg - \mu v}{m}.
\end{cases}
\]

Начальные условия:
\[
h(0) = h_0, \quad v(0) = 0.
\]

Сила натяжения троса $T(t)$ является управляющим параметром и зависит от режима работы лифта. Рассмотрим основные режимы:

\textbf{Режим разгона (подъем):}
При разгоне лифта сила натяжения максимальна:
\[
T(t) = T_{\max}, \quad \text{при } v < v_{\max} \text{ и } a < a_{\max}.
\]

В этом режиме ускорение определяется как:
\[
a(t) = \frac{T_{\max} - mg - \mu v}{m}.
\]

\textbf{Режим равномерного движения:}
При достижении максимальной скорости сила натяжения компенсирует силу тяжести и трения:
\[
T(t) = mg + \mu v_{\max}.
\]

В этом режиме ускорение равно нулю, скорость постоянна:
\[
a(t) = 0, \quad v(t) = v_{\max} = \text{const}.
\]

\textbf{Режим торможения:}
При торможении сила натяжения уменьшается:
\[
T(t) = mg - ma_{\max} + \mu v.
\]

Ускорение в режиме торможения:
\[
a(t) = \frac{T(t) - mg - \mu v}{m} = -a_{\max}.
\]

\textbf{Режим спуска:}
При спуске сила натяжения минимальна:
\[
T(t) = mg - ma_{\max} - \mu v.
\]

Ускорение при спуске:
\[
a(t) = \frac{T(t) - mg - \mu v}{m} = -a_{\max} - \frac{\mu v}{m}.
\]

\textbf{Аналитическое решение для режима разгона без трения:}
При $\mu = 0$ и постоянной силе $T = T_{\max}$ система имеет аналитическое решение:
\[
v(t) = \frac{T_{\max} - mg}{m}t, \quad h(t) = h_0 + \frac{T_{\max} - mg}{2m}t^2.
\]

Время достижения максимальной скорости:
\[
t_1 = \frac{m v_{\max}}{T_{\max} - mg}.
\]

\textbf{Аналитическое решение с учетом трения:}
При $\mu \ne 0$ и постоянной силе $T = T_{\max}$ решение для скорости имеет вид:
\[
v(t) = \frac{T_{\max} - mg}{\mu}\left(1 - e^{-\frac{\mu}{m}t}\right).
\]

Высота определяется интегрированием скорости:
\[
h(t) = h_0 + \frac{T_{\max} - mg}{\mu}\left(t - \frac{m}{\mu}\left(1 - e^{-\frac{\mu}{m}t}\right)\right).
\]

Установившаяся скорость (при $t \to \infty$):
\[
v_{\infty} = \frac{T_{\max} - mg}{\mu}.
\]

\section{Постановка и решение вычислительных задач}
После формализации математической модели необходимо поставить и решить набор соответствующих вычислительных задач для получения численных результатов.

\subsection{Постановка задачи Коши}
Система ОДУ первого порядка с начальными условиями представляет собой корректную задачу Коши:
\[
\begin{cases}
  \dfrac{dh}{dt} = v,\\
  \dfrac{dv}{dt} = \dfrac{T(t) - mg - \mu v}{m},\\
  h(0) = h_0, \quad v(0) = 0.
\end{cases}
\]

Правая часть системы непрерывна и липшицева по переменным состояния, что гарантирует существование и единственность решения.

\subsection{Выбор временного интервала}
Временной интервал выбирается в зависимости от режима движения:
\begin{itemize}
  \item Для режима разгона: $t \in [0, t_1]$, где $t_1$ — время достижения максимальной скорости
  \item Для режима равномерного движения: $t \in [t_1, t_2]$, где $t_2$ — начало торможения
  \item Для режима торможения: $t \in [t_2, t_3]$, где $t_3$ — момент остановки
\end{itemize}

Время разгона до максимальной скорости:
\[
t_1 = \frac{v_{\max}}{a_{\max}}.
\]

Время торможения:
\[
t_3 - t_2 = \frac{v_{\max}}{a_{\max}}.
\]

\subsection{Численное интегрирование}
Для численного решения системы дифференциальных уравнений используется метод Рунге--Кутты 4-го порядка (RK4). Этот метод обеспечивает высокую точность и устойчивость решения.

Система уравнений в векторной форме:
\[
\frac{dy}{dt} = f(t, y),
\]
где
\[
y = \begin{pmatrix} h \\ v \end{pmatrix}, \quad f(t, y) = \begin{pmatrix} v \\ \dfrac{T(t) - mg - \mu v}{m} \end{pmatrix}.
\]

Формулы метода Рунге--Кутты 4-го порядка:
\[
\begin{aligned}
k_1 &= \Delta t \cdot f(t_n, y_n),\\
k_2 &= \Delta t \cdot f\!\left(t_n+\frac{\Delta t}{2},\, y_n+\frac{k_1}{2}\right),\\
k_3 &= \Delta t \cdot f\!\left(t_n+\frac{\Delta t}{2},\, y_n+\frac{k_2}{2}\right),\\
k_4 &= \Delta t \cdot f(t_n+\Delta t,\, y_n+k_3),\\
y_{n+1} &= y_n+\frac{k_1+2k_2+2k_3+k_4}{6}.
\end{aligned}
\]

Распишем вычисления для каждого шага подробно. Пусть на шаге $n$ известны значения $h_n$ и $v_n$.

\textbf{Шаг 1: вычисление $k_1$}
\[
k_1 = \Delta t \cdot \begin{pmatrix} v_n \\ \dfrac{T(t_n) - mg - \mu v_n}{m} \end{pmatrix} = \begin{pmatrix} \Delta t \cdot v_n \\ \Delta t \cdot \dfrac{T(t_n) - mg - \mu v_n}{m} \end{pmatrix}.
\]

Обозначим компоненты: $k_1 = [k_{1h}, k_{1v}]^T$, тогда:
\[
k_{1h} = \Delta t \cdot v_n, \quad k_{1v} = \Delta t \cdot \frac{T(t_n) - mg - \mu v_n}{m}.
\]

\textbf{Шаг 2: вычисление $k_2$}
\[
k_2 = \Delta t \cdot f\!\left(t_n+\frac{\Delta t}{2},\, y_n+\frac{k_1}{2}\right) = \Delta t \cdot \begin{pmatrix} v_n + \frac{k_{1v}}{2} \\ \dfrac{T(t_n+\Delta t/2) - mg - \mu(v_n + k_{1v}/2)}{m} \end{pmatrix}.
\]

Компоненты:
\[
k_{2h} = \Delta t \cdot \left(v_n + \frac{k_{1v}}{2}\right), \quad k_{2v} = \Delta t \cdot \frac{T(t_n+\Delta t/2) - mg - \mu(v_n + k_{1v}/2)}{m}.
\]

\textbf{Шаг 3: вычисление $k_3$}
\[
k_3 = \Delta t \cdot f\!\left(t_n+\frac{\Delta t}{2},\, y_n+\frac{k_2}{2}\right) = \Delta t \cdot \begin{pmatrix} v_n + \frac{k_{2v}}{2} \\ \dfrac{T(t_n+\Delta t/2) - mg - \mu(v_n + k_{2v}/2)}{m} \end{pmatrix}.
\]

Компоненты:
\[
k_{3h} = \Delta t \cdot \left(v_n + \frac{k_{2v}}{2}\right), \quad k_{3v} = \Delta t \cdot \frac{T(t_n+\Delta t/2) - mg - \mu(v_n + k_{2v}/2)}{m}.
\]

\textbf{Шаг 4: вычисление $k_4$}
\[
k_4 = \Delta t \cdot f(t_n+\Delta t,\, y_n+k_3) = \Delta t \cdot \begin{pmatrix} v_n + k_{3v} \\ \dfrac{T(t_n+\Delta t) - mg - \mu(v_n + k_{3v})}{m} \end{pmatrix}.
\]

Компоненты:
\[
k_{4h} = \Delta t \cdot (v_n + k_{3v}), \quad k_{4v} = \Delta t \cdot \frac{T(t_n+\Delta t) - mg - \mu(v_n + k_{3v})}{m}.
\]

\textbf{Шаг 5: вычисление нового значения}
\[
y_{n+1} = y_n + \frac{k_1 + 2k_2 + 2k_3 + k_4}{6},
\]

что в компонентах дает:
\[
\begin{aligned}
h_{n+1} &= h_n + \frac{k_{1h} + 2k_{2h} + 2k_{3h} + k_{4h}}{6},\\
v_{n+1} &= v_n + \frac{k_{1v} + 2k_{2v} + 2k_{3v} + k_{4v}}{6}.
\end{aligned}
\]

\textbf{Пример вычисления одного шага:}
Пусть $m = 1000$ кг, $g = 9.81$ м/с², $\mu = 50$ Н·с/м, $T(t) = 15000$ Н (режим разгона), $\Delta t = 0.1$ с, $h_0 = 0$ м, $v_0 = 0$ м/с.

На первом шаге ($n = 0$):
\[
k_{1h} = 0.1 \cdot 0 = 0, \quad k_{1v} = 0.1 \cdot \frac{15000 - 1000 \cdot 9.81 - 50 \cdot 0}{1000} = 0.1 \cdot 5.19 = 0.519.
\]

\[
k_{2h} = 0.1 \cdot \left(0 + \frac{0.519}{2}\right) = 0.02595,
\]
\[
k_{2v} = 0.1 \cdot \frac{15000 - 9810 - 50 \cdot 0.2595}{1000} = 0.1 \cdot \frac{5187.025}{1000} = 0.5187.
\]

\[
k_{3h} = 0.1 \cdot \left(0 + \frac{0.5187}{2}\right) = 0.025935,
\]
\[
k_{3v} = 0.1 \cdot \frac{15000 - 9810 - 50 \cdot 0.25935}{1000} = 0.1 \cdot \frac{5187.0325}{1000} = 0.5187.
\]

\[
k_{4h} = 0.1 \cdot (0 + 0.5187) = 0.05187,
\]
\[
k_{4v} = 0.1 \cdot \frac{15000 - 9810 - 50 \cdot 0.5187}{1000} = 0.1 \cdot \frac{5164.065}{1000} = 0.5164.
\]

Новые значения:
\[
h_1 = 0 + \frac{0 + 2 \cdot 0.02595 + 2 \cdot 0.025935 + 0.05187}{6} = 0.0259 \text{ м},
\]
\[
v_1 = 0 + \frac{0.519 + 2 \cdot 0.5187 + 2 \cdot 0.5187 + 0.5164}{6} = 0.5183 \text{ м/с}.
\]

Погрешность метода Рунге--Кутты 4-го порядка составляет $O(\Delta t^5)$ на каждом шаге, что обеспечивает высокую точность численного решения.

\subsection{Критерии переключения режимов}
Переключение между режимами происходит при выполнении условий:
\begin{itemize}
  \item Переход от разгона к равномерному движению: $v(t) \ge v_{\max}$
  \item Начало торможения: $h(t) \ge H - \frac{v_{\max}^2}{2a_{\max}}$
  \item Остановка: $v(t) \le \varepsilon$ и $|h(t) - H| \le \delta$, где $\varepsilon$ и $\delta$ — малые положительные числа
\end{itemize}

\section{Анализ моделей и проверка качества}
После построения математической модели и решения вычислительных задач необходимо проверить качество модели, проанализировать полученные результаты и оценить применимость модели для описания реального процесса.

\subsection{Анализ решений}
Математическая модель лифта позволяет проанализировать различные аспекты движения кабины.

При отсутствии трения ($\mu = 0$) и постоянной силе натяжения $T = \text{const}$ система имеет аналитическое решение:
\[
h(t) = h_0 + \frac{T - mg}{2m}t^2, \quad v(t) = \frac{T - mg}{m}t.
\]

Учет трения приводит к экспоненциальному установлению скорости:
\[
v(t) = \frac{T - mg}{\mu}\left(1 - e^{-\frac{\mu}{m}t}\right).
\]

Влияние параметров на движение:
\begin{itemize}
  \item Увеличение массы $m$ приводит к уменьшению ускорения при той же силе натяжения
  \item Увеличение коэффициента трения $\mu$ уменьшает максимальную скорость и увеличивает время разгона
  \item Максимальная сила натяжения $T_{\max}$ определяет максимальное ускорение лифта
\end{itemize}

\subsection{Проверка качества модели}
Для проверки качества модели необходимо сравнить результаты, полученные с помощью математической модели, с реальным поведением системы или с эталонными решениями.

Критерии качества модели:
\begin{itemize}
  \item \textbf{Корректность формализации:} система дифференциальных уравнений правильно описывает физические законы движения
  \item \textbf{Адекватность допущений:} упрощения модели не приводят к существенным ошибкам в описании процесса
  \item \textbf{Точность численного решения:} метод интегрирования обеспечивает достаточную точность для практических задач
  \item \textbf{Устойчивость решения:} малые изменения параметров не приводят к качественным изменениям поведения системы
\end{itemize}

Ограничения модели: не учитываются упругие деформации троса, инерция двигателя, влияние ветра и вибрации, изменение массы при входе/выходе пассажиров. Эти факторы могут влиять на точность модели в реальных условиях, что необходимо учитывать при практическом применении.

\textbf{Качественный анализ поведения системы:}

При разгоне лифта (режим 1) скорость увеличивается по закону:
\[
v(t) = \frac{T_{\max} - mg}{\mu}\left(1 - e^{-\frac{\mu}{m}t}\right),
\]
что приводит к экспоненциальному приближению к установившейся скорости $v_{\infty} = (T_{\max} - mg)/\mu$.

В режиме равномерного движения скорость постоянна, высота изменяется линейно:
\[
h(t) = h(t_1) + v_{\max}(t - t_1),
\]
где $t_1$ --- момент начала равномерного движения.

При торможении скорость уменьшается по закону:
\[
v(t) = v_{\max} - a_{\max}(t - t_2),
\]
где $t_2$ --- момент начала торможения. Высота при торможении:
\[
h(t) = h(t_2) + v_{\max}(t - t_2) - \frac{a_{\max}}{2}(t - t_2)^2.
\]

\section{Результаты}
В результате построения математической модели и решения вычислительных задач получены следующие результаты:

\begin{itemize}
  \item \textbf{Формализация процесса:} движение лифта формализовано через систему дифференциальных уравнений первого порядка, описывающую динамику высоты и скорости кабины
  \item \textbf{Модель управления:} определены режимы работы лифта (разгон, равномерное движение, торможение, спуск) и соответствующие им значения силы натяжения троса
  \item \textbf{Критерии переключения режимов:} установлены условия перехода между режимами на основе текущих значений высоты и скорости
  \item \textbf{Аналитические решения:} получены аналитические выражения для частных случаев (без трения, постоянная сила), которые могут служить для проверки качества численного решения
  \item \textbf{Численный метод:} предложен метод Рунге-Кутты 4-го порядка для решения системы дифференциальных уравнений
\end{itemize}

\section{Выводы}
\begin{itemize}
  \item Выполнена формализация процесса движения лифта на языке математических формул. Процесс движения кабины описан через систему дифференциальных уравнений первого порядка, что является основой математической модели.
  
  \item Построение математической модели включает описание физических законов (второй закон Ньютона), введение переменных состояния (высота и скорость) и формализацию управляющих воздействий (сила натяжения троса в различных режимах).
  
  \item Поставлены и решены вычислительные задачи: задача Коши для системы дифференциальных уравнений, определение критериев переключения режимов, выбор метода численного интегрирования (метод Рунге-Кутты 4-го порядка).
  
  \item Проведена проверка качества модели: получены аналитические решения для частных случаев, проанализировано влияние параметров системы на характер движения, выявлены ограничения модели, связанные с принятыми допущениями.
  
  \item Модель учитывает основные физические факторы: силу тяжести, силу натяжения троса и силу трения, пропорциональную скорости движения. Определены четыре основных режима работы лифта: разгон, равномерное движение, торможение и спуск.
  
  \item Показано, что учет трения приводит к экспоненциальному установлению скорости и влияет на время разгона и максимальную скорость движения. Модель может быть использована для анализа влияния параметров системы на характеристики движения лифта.
  
  \item Ограничения модели связаны с упрощениями: не учитываются упругие деформации троса, инерция двигателя, динамическое изменение массы и другие факторы. Для повышения качества модели необходимо учитывать дополнительные физические эффекты или проводить натурные эксперименты для верификации результатов.
\end{itemize}

\begin{thebibliography}{9}
\bibitem{runge_kutta}
Бахвалов Н.С., Жидков Н.П., Кобельков Г.М. Численные методы. --- М.: БИНОМ. Лаборатория знаний, 2011. --- 636 с.

\bibitem{differential_equations}
Понтрягин Л.С. Обыкновенные дифференциальные уравнения. --- М.: Наука, 1974. --- 331 с.

\bibitem{mechanics}
Голубев Ю.Ф. Основы теоретической механики. --- М.: Изд-во МГУ, 2000. --- 720 с.
\end{thebibliography}

\end{document}

